\documentclass{jsarticle}
\usepackage[dvipdfmx, final]{graphicx}
\usepackage{amsmath}
\usepackage{amssymb}
\usepackage{chemfig}
\begin{document}
\title{}
\author{}
\date{}
\maketitle

   鈴木杞憂という、彩さんの影武者の人が、この人は、2014年10月4日の
   午前6時30分に、福山駅の待合の椅子に、プロメトリックの参考書を
   引っさげて、私が、初対面だったので、何？をしているかを、
   聞きづらく、「どんなことを、勉強しているのですか？」と
   聞くと、左に置いている参考書の鞄を持とうとしていたので、
   去るのかーと、「私が、去ります」で、私が去っていった。
   彩さんの嫉妬心を、、、

   街灯の下に、ある白い口髭を生やした年老いた人が、
   アインシュタインと見えて、
   手招きして、手を振りながら、私に近づいて、
   今まで、見ていなかったように、突然姿が消えて、
   姉が、「この世では、秘密のドアがあり、
   そのネットショップから、お金がもらえる、ドアがあるんだよ。」
   と言われて、ある深夜の、今何時と、00時00分と、スマフォの光
   が見えて、そしたら、深い夢の中で、ドン、ドン、ドン、ドン、
   とどこからか、音が響いて、どんどん、音が、大きくなり、
   横の部屋は、もう誰も住んでいないし、上の人かな？と、
   布団に潜り、身を縮めて、そしたら、スマフォの時間が。、
   00時33分であり、夢の中で、
   街灯の下に、ある白い口髭を生やした年老いた人が、
   アインシュタインと見えて、
   手招きして、手を振りながら、私に近づいて、
   今まで、見ていなかったように、突然姿が消えて、
   耳元で、ほとんど破れた年老いた老人の声が、耳元で、
   「やっと、みつけた」
   と、すると、側の影の姿が、別の人に重なり、融合して、

   この街には、嘘をつくと、鼻が長くなるという、
   言い伝えがあり、思っていることと、違うことを
   思って、言っていたら、、、

   アインシュタイン博士と、タイプの女性が同じらしく、

   この後に、ネットショップをしていたら、お金がどんどん、
   入り、700万円儲かり、
   こんな話、マサ君、アインシュタインに会ったことないの？

  私は、2013年の1月29日の深夜、2時00分から3時00分にかけて、
  1月の29日から1月の31日の、2時から3時にかけて、
  アカシックレコードにアクセスしようと思ったきっかけの犯人は、
  宇多田ヒカルさんで、してみると、情報空間の知らない経路を
  通っていたら、ゴミと一緒に、$x\log x$という、式をアインシュタイン
  博士に、自分で、この人にアクセスしようと思っていたので、
  渡されて、これ何？と聞くと、ゼータ関数だよと、
  その後に、調べまくっていたら、非可換代数方程式
  $\pi = [i\pi(x),f(x)]$という式がネットショップから出てきて、
  これも調べていたら、$\pi = \int x\log xdx$という式も出てきて、
  それを、ポワンカレ予想を解いた人の論文を見ていたら、
  ${d \over {d f}}F(x,y) \ge 0$までも、出てきて、
  それから、${d \over {d f}}F(x,y) = \log (x \log x) \ge {2\sqrt{(y \log
  y)}}$も、出てきて、
  私の子どもか分からない子に、アカシックレコードのネットショップ
  にアクセスしたら、$H\Psi = \bigoplus{({i\hbar}^{\nabla})}^{\oplus L}$
  という式をもらって、調べると、${d \over {d f}}F(x,y) = F^{f^{'}}$
  と、$\int \Gammma^{'}(\gamma)dx_m= \Gamma^{\gamma}$という、
  式も分かり、
  彩さんの影武者から聞いた、ネットショップって、アカシックレコード
  なんだと、ふ〜んと、
  アインシュタイン博士の隣に、月読様の盧遮那仏様が、居て、
  アインシュタイン博士は、姿を消していて、月読様が居るのが、
  神様と見えて、月読様は、相手によって、姿を変えるので、


  　

\end{document}


